\section*{Bab \ref{ch:g7:prop} --- Perbandingan Senilai}

\begin{solution}{exo:g7:prop:1}
Tabel pertama: \omterm{def:g3:division:remainder}{hasil baginya} $\frac{7.5}{3} = 2.5$,
$\frac{17.5}{7} = 2.5$, $\frac{27.5}{11} = 2.5$: \omterm{def:g6:propdata:def}{sebanding}, dengan
koefisien $2.5$.

Tabel kedua: $\frac62 = 3$ dan $\frac{15}{5} = 3$, tetapi
$\frac{28}{9} \approx 3.1$: tidak \omterm{def:g6:propdata:def}{sebanding}.
\end{solution}

\begin{solution}{exo:g7:prop:2}
Dengan aturan perkalian silang: $5 \times p = 6 \times 8$, jadi
$p = \frac{48}{5} = 9.60$ euro.
\end{solution}

\begin{solution}{exo:g7:prop:3}
Lajunya: $36 \div 3 = 12$ halaman per menit. Dalam $5$ menit:
$60$ halaman. Untuk $96$ halaman: $96 \div 12 = 8$ menit.
\end{solution}

\begin{solution}{exo:g7:prop:4}
$30\,\%$ dari $250$: $0.3 \times 250 = 75$. \quad
$15\,\%$ dari $60$: $0.15 \times 60 = 9$. \quad
$t\,\%$ dari $80$ adalah $20$: $\frac{20}{80} = \frac{25}{100}$, jadi
$t = 25$.
\end{solution}

\begin{solution}{exo:g7:prop:5}
$\frac{18}{25} = \frac{18 \times 4}{25 \times 4} = \frac{72}{100} =
72\,\%$.
\end{solution}

\begin{solution}{exo:g7:prop:6}
Jalan setapaknya: $9 \times 25\,000 = 225\,000$ cm $= 2.25$ km.
Danaunya: $2$ km $= 200\,000$ cm di tanah, jadi
$200\,000 \div 25\,000 = 8$ cm pada peta.
\end{solution}

\begin{solution}{exo:g7:prop:7}
$4.3$ m $= 430$ cm, dan $430 \div 43 = 10$ cm.
\end{solution}

\begin{solution}{exo:g7:prop:8}
Cokelat per sajian: $240 \div 8 = 30$ g. Dengan $300$ g:
$300 \div 30 = 10$ sajian tepat.
\end{solution}

\begin{solution}{exo:g7:prop:9}
\emph{1.} Potongannya: $20$ euro, yaitu $\frac{20}{80} = 25\,\%$ dari
harga semula.

\emph{2.} Kenaikan $20$ euro itu sekarang dibandingkan dengan acuan
yang \emph{baru}, yaitu $60$: $\frac{20}{60} = \frac13 \approx 33\,\%$.
Persen selalu mengacu pada suatu keseluruhan; keseluruhannya berubah,
jadi persennya berubah juga.
\end{solution}

\begin{solution}{exo:g7:prop:10}
Koefisien dari kolom yang lengkap: $\frac{17.5}{10} = 1.75$. Lalu
$y = 4 \times 1.75 = 7$ dan $x = \frac{10.5}{1.75} = 6$.
\end{solution}

\begin{solution}{exo:g7:prop:11}
\emph{1.} Setelah satu salinan: $10 \times 0.8 = 8$ cm. Setelah dua:
$8 \times 0.8 = 6.4$ cm.

\emph{2.} Pengecilan yang kedua berlaku pada salinan yang sudah
diperkecil, bukan pada aslinya: faktornya saling dikalikan,
$0.8 \times 0.8 = 0.64$, jadi dua salinan memperkecilnya menjadi
$64\,\%$ ukuran aslinya --- bukan $60\,\%$.
\end{solution}

\begin{solution}{pb:g7:prop:1}
\textbf{1.} Tinggi dan bayangan saling \omterm{def:g6:propdata:def}{sebanding}, jadi menurut
aturan perkalian silang (\cref{thm:g7:prop:cross}), dengan kolom
tongkatnya $(1,\ 0.4)$ dan kolom pohonnya $(h,\ 3.2)$:
$0.4 \times h = 1 \times 3.2$, sehingga
$h = \frac{3.2}{0.4} = 8$ m.

\textbf{2.} Menaranya: $\frac{26}{0.4} = 65$ m. Bayangan anaknya:
$1.5 \times 0.4 = 0.6$ m (bayangan $=$ tinggi $\times$ koefisien
saat itu, $0.4$).

\textbf{3.} Ketika matahari bergerak melintasi langit, koefisien
yang menghubungkan tinggi dengan bayangan ikut berubah; hanya
pengukuran pada saat yang sama yang berbagi koefisien yang sama.

\textbf{4.} Pada saat itu koefisiennya $1$: setiap tinggi sama dengan
bayangannya. Puncak piramidanya karena itu setinggi
$147$ m --- yaitu tinggi Piramida Besar (ujung bayangannya, yang
diukur dari titik di bawah puncaknya, sudah cukup).

\textbf{5.} \omterm{def:g3:division:remainder}{Hasil bagi} kolomnya
$\frac{\text{bayangan}}{\text{tinggi}}$ --- yaitu koefisien
\omterm{def:g7:prop:table}{perbandingan senilai} pada saat itu --- sama bagi setiap benda yang
tegak: persis itulah yang membuat tabel tinggi dan bayangan menjadi
\omterm{def:g7:prop:table}{tabel perbandingan senilai} (\cref{def:g7:prop:table}).

\textbf{6.} Sinar matahari datang \omterm{def:g6:lines:perp}{sejajar}. Di Bumi yang \emph{datar},
dua sinar \omterm{def:g6:lines:perp}{sejajar} akan mengenai dua tongkat tegak dengan sudut yang
sama: keduanya akan membuat bayangan yang \omterm{def:g6:propdata:def}{sebanding} --- entah
keduanya tanpa bayangan, entah keduanya berbayangan. Satu tongkat
tanpa bayangan (Syene) dan satu lagi berbayangan (Aleksandria), pada
saat yang sama, mustahil di Bumi yang datar: kedua arah tegaknya
pasti menunjuk arah yang berbeda --- permukaannya melengkung.

\textbf{7.} Pada gambarnya, arah tegak di Syene menunjuk lurus ke
matahari; arah tegak di Aleksandria miring sebesar sudut antara arah
kedua kota itu, dilihat dari pusatnya. Karena sinarnya \omterm{def:g6:lines:perp}{sejajar},
kemiringan itu tepat sebesar $7.2^\circ$ yang terukur antara sinar
dan tongkat di Aleksandria.

\textbf{8.} $\frac{7.2}{360} = \frac{72}{3600} = \frac{1}{50}$:
seperlima puluh satu putaran penuh.

\textbf{9.} Panjang busur \omterm{def:g6:propdata:def}{sebanding} dengan sudutnya: kalau
$7.2^\circ$ --- seperlima puluh putaran --- bersesuaian dengan
$800$ km, maka satu putaran penuh bersesuaian dengan
\[
50 \times 800 = 40\,000 \text{ km}.
\]

\textbf{10.} \omterm{def:g4:geometry:circle}{Diameter} $= \frac{\text{keliling}}{\pi}
\approx \frac{40\,000}{3.14} \approx 12\,700$ km --- melawan
$12\,742$ km yang modern: benar sampai jauh di bawah satu persen,
hanya dengan sebatang tongkat.

\textbf{11.} $40\,000$ km $= 4\,000\,000\,000$ cm, dan
$4\,000\,000\,000 \div 40\,000\,000 = 100$ cm $= 1$ m
\omterm{def:g6:measure:perimeter}{kelilingnya}. \omterm{def:g4:geometry:circle}{Diameternya}: $100 \div 3.14 \approx 32$ cm --- sebuah
bola dunia meja yang gagah.

\textbf{12.} $4.8$ km $= 480\,000$ cm, dan
$480\,000 \div 40\,000\,000 = 0.012$ cm $= 0.12$ mm. Gunung
tertinggi di Pegunungan Alpen hanya sepersepuluh milimeter pada bola
dunia yang \omterm{def:g6:measure:perimeter}{kelilingnya} semeter: sebutir debu. Pada \omterm{def:g7:prop:scale}{skala} ini Bumi
lebih halus daripada kebanyakan bola yang dipoles.

\textbf{13.} $40\,000 \div 40 = 1\,000$ hari, dan
$1\,000 \div 365 \approx 2.7$: hampir tiga tahun berjalan kaki.

\textbf{14.} $\frac{10}{6\,370} \approx 0.0016$, yaitu kira-kira
$0.2\,\%$ (lebih tepatnya $0.16\,\%$): atmosfernya adalah kulit
setipis bisikan pada planet ini.

\textbf{15.} Kekeliruannya: $42\,000 - 40\,000 = 2\,000$ km, dan
$\frac{2\,000}{40\,000} = \frac{5}{100} = 5\,\%$. Satu kalimat:
dengan sebatang tongkat, sebuah sumur, sebuah jalan yang terukur dan
satu \omterm{def:g7:prop:table}{perbandingan senilai}, Eratosthenes mengukur sebuah planet sampai
\omterm{ex:g1:subtraction:difference}{selisih} beberapa persen --- matematika melaju jauh dengan bekal yang
sangat sedikit.

\end{solution}
